% !TeX root = ../../../navier-stokes-manuscript.tex
% ============================================================
% 2.2 Finite-Time Response Escalation
% ============================================================

\subsection{Finite-Time Response Escalation}\label{sub:FiniteTimeResponseEscalation}

Follow a particle path \(X(t)\) inside the classical interval, so that \(\dot X(t)=u(X(t),t)\).
Along that path, differentiation is the material derivative \(D_t=\partial_t+u\cdot\nabla\).
For an integer \(r\geq0\), write the \(r\)-th carried response as
\[
  A_r(t)
  :=
  D_t^r u(X(t),t).
\]
Classical smoothness before \(T_*\) gives
\[
  \frac{d}{dt}A_r(t)=A_{r+1}(t).
\]
The fundamental theorem of calculus gives, for \(t_0<t<T_*\),
\[
  A_r(t)
  =
  A_r(t_0)
  +
  \int_{t_0}^{t}A_{r+1}(\tau)\,d\tau.
\]
Taking the Euclidean norm yields
\[
  |A_r(t)|
  \leq
  |A_r(t_0)|
  +(T_*-t_0)
  \sup_{t_0\leq\tau<T_*}|A_{r+1}(\tau)|.
\]
If the next response stays uniformly bounded, then the present response stays uniformly bounded across the remaining finite time.
Its contrapositive is the useful direction:
\[
  \sup_{t_0\leq t<T_*}|A_r(t)|=\infty
  \quad\Longrightarrow\quad
  \sup_{t_0\leq t<T_*}|A_{r+1}(t)|=\infty.
\]
Applying the same implication again shows that an unbounded rung cannot sit beneath a bounded finite ceiling of higher material responses.

This statement concerns a response rung that actually becomes unbounded.
It says nothing about the sign of a response, and it gives no monotonicity.
A smooth rung may vanish, reverse direction, or oscillate many times.
Steady solutions and the zero solution remain fully compatible with the argument because none of their rungs is assumed to blow up.
The implication becomes active only inside the Blown branch whose finite-time escape we are trying to understand.

The particle calculation supplies the physical meaning, while estimates compare the field at fixed spatial coordinates.
The same FTOC mechanism works in any Banach space \(X\).
If \(Z\in C^1([t_0,T_*);X)\), then
\[
  \norm{Z(t)}_X
  \leq
  \norm{Z(t_0)}_X
  +
  \int_{t_0}^{t}\norm{Z'(\tau)}_X\,d\tau.
\]
In particular, a uniform bound for \(Z'\) on a finite interval forces a uniform bound for \(Z\).
We may take
\[
  Z(t)
  =
  \partial_t^m\partial_x^\alpha u(\,cdot\,,t)
\]
in any Sobolev space where the classical solution makes the expression meaningful.
Then \(Z'(t)=\partial_t^{m+1}\partial_x^\alpha u(\,cdot\,,t)\), so the same temporal escalation is attached to every spatial derivative \(\partial_x^\alpha u\).

This is the point at which one derivative index becomes two.
The integer \(m\) records temporal response at a fixed spatial coordinate, while the multi-index \(\alpha\) records spatial response inside the field.
Neither index determines the other.
A failure could first become visible at one temporal depth and one spatial height, so both have to remain available independently.

The material responses are built from these fixed-coordinate derivatives.
For example,
\[
  D_tu
  =
  \partial_tu+(u\cdot\nabla)u,
\]
and one more carried derivative gives
\begin{align*}
  D_t^2u
  ={}&
  \partial_t^2u
  +2(u\cdot\nabla)\partial_tu
  +(\partial_tu\cdot\nabla)u
  \\
  &+
  (u\cdot\nabla)\bigl((u\cdot\nabla)u\bigr).
\end{align*}
The jerk already contains pure time derivatives, spatial derivatives, and products between lower responses.
Higher material responses continue this proliferation.
The fixed-coordinate family \(\{\partial_t^m\partial_x^\alpha u\}\) is the natural set of coordinates from which the physical response tower is assembled.

These coordinates have not been imposed on the fluid from outside.
They are already nested in the original equation, and direct differentiation reveals exactly how they relate.
